\clearpage
\fancyhead[L]{\small PALLAS. JUNO.}
\hypertarget{native-274-lower}{}
\begin{center}
PALLAS. JUNO.

\rule{0.6\linewidth}{0.4pt}

Monatliche Correspondenz. Band XIII. S. 313–315. März 1806.

\rule{0.6\linewidth}{0.4pt}
\end{center}

\par\indent Von \emph{Pallas} sind uns bisher nur die drei Beobachtungen bekannt geworden, die an der Seeberger Sternwarte gemacht und im vorigen Heft auf \href{https://www.digitale-sammlungen.de/en/view/bsb10538605?page=197}{S. 189} mitgeteilt wurden. Von \emph{Juno} ist uns hingegen noch keine einzige bekannt. Deshalb beeilen wir uns, unseren Lesern die folgende Nachricht mitzuteilen.

\par\indent Dr.\ \textsc{Gauss} sah \emph{Pallas} am 13. Februar zum ersten Mal wieder. Er machte diese Beobachtung nur einmal und allein zu dem Zweck, den Planeten am nächsten klaren Abend anhand seiner Bewegung unter den anderen benachbarten Sternen herauszufinden. Wegen der Unsicherheit musste er sie alle ebenfalls festhalten. Weil er die Beobachtung für sehr ungenau hielt, schickte er sie uns gar nicht. Am folgenden Abend, dem 14. Februar, fiel es ihm dann leicht, die weitergezogene Pallas zu erkennen, die wie ein Stern 7\textsuperscript{ter} bis 8\textsuperscript{ter} Größe erschien. An diesem Abend verglich er sie mit 54 Eridani, an den folgenden Abenden mit Sternen der Hist.\ cél. Hier sind die Beobachtungen [dieselben wie auf \hyperlink{native-272}{S. 267} dieses Bandes].

\par\indent Diesen Beobachtungen zufolge gibt Dr.\ \textsc{Gauss}’ Ephemeride in der M.\,C., Band XI, \href{https://www.digitale-sammlungen.de/en/view/bsb10538603?page=388}{S. 376} [\hyperlink{pallas180504.1}{S. 263} dieses Bandes] die Rektaszension dieses Planeten um 27 Minuten und die Deklination um 5 Minuten zu groß an. „Es ist sonderbar“, bemerkt Dr.\ \textsc{Gauss}, „dass die erste Abweichung genau in die entgegengesetzte Richtung weist wie die, die sich aus Ihren Beobachtungen vom November des vorigen Jahres ergibt (siehe das vorige Heft, \href{https://www.digitale-sammlungen.de/en/view/bsb10538605?page=197}{S. 189}). In der Größe des Fehlers zeigt sich jetzt, wie ich glaube, der Einfluss der Störungen.“

\par\indent \emph{Juno} sah er am 16. Februar zum ersten Mal wieder, und am 17. Februar beobachtete er sie wieder als \emph{solche}. Hier sind seine Beobachtungen [dieselben wie auf \hyperlink{native-272}{S. 267} dieses Bandes].

\par\indent Danach würde Dr.\ \textsc{Gauss}’ Ephemeride (M.\,C., Band XI, \href{https://www.digitale-sammlungen.de/en/view/bsb10538603?page=491}{S. 477} [\hyperlink{native-270}{S. 265} dieses Bandes]) die Rektaszension um eine oder anderthalb Minuten zu groß und die Deklination um eine Minute zu klein angeben. Nach genau einem Jahr ist das eine ganz unbedeutende Abweichung. Daraus lässt sich zugleich schließen, dass Dr.\ \textsc{Gauss}’ letzte Beobachtung vom 20. Februar 1805, auf der seine V.\ Ele%
\SourceLeaf{275}%
mente beruhten, besonders gut gewesen sein muss. Dr.~\textsc{Gauss} schreibt bei dieser Gelegenheit: „Juno ist sehr lichtschwach. Ich schätze sie etwa so hell wie einen Stern elfter Größe. Bei klarem Himmel lässt sie sich mit dem Kreismikrometer noch gut beobachten. Mit dem Passageninstrument und dem Kreis oder Quadranten wird das aber kaum möglich sein, wenn Beleuchtung nötig ist. Diese Beobachtungen sind für mich gewöhnlich sehr mühsam, weil ich dazu nur ein Spiegelteleskop von \textsc{Short} verwenden kann. Sein Gesichtsfeld ist sehr klein. Daher fehlen meistens geeignete Sterne mit gut bestimmter Position, die nahe genug liegen und zusammen mit dem Planeten durch das enge Feld ziehen. Hinzu kommt die oft erhebliche Schwierigkeit, solche kleinen Planeten überhaupt erst zu finden, wenn keine bekannten Sterne nahe genug sind, um als Leitsterne zu dienen.“

\par\indent Diesen Nachrichten fügen wir eine weitere hinzu, die unseren astronomisch interessierten Lesern gewiss höchst willkommen sein wird. Dr.~\textsc{Gauss} berichtet uns, dass er sich jetzt hauptsächlich mit verschiedenen Einzelfragen seiner Methode zur Bestimmung der Planetenbahnen beschäftigt. Er hofft, bald ernsthaft mit der Ausarbeitung, Fertigstellung und Herausgabe eines eigenen Werks darüber beginnen zu können.

\begin{flushright}
[\textsc{von Zach}.]
\end{flushright}
